%! Periodic Phenomena — Unit 3, Lesson 3.1 — Warm-Up (Answer Key)
\documentclass[10pt]{article}
\usepackage{precalculus-article}
\usepackage{precalculus-key}   % pulls in -boxes; do NOT also load -boxes
\newcommand{\TallMath}[1]{$\displaystyle #1\rule[-1.4em]{0pt}{3.2em}$}

\begin{document}

\pageheader{Unit 3, Lesson 3.1}{Warm-Up --- Answer Key}
\namedateperiod

\vspace{0.12in}

\begin{notesbox}{Spiral Review --- Key}

\textbf{1.}\quad The table below shows values of a function $f$.

\vspace{4pt}
\begin{center}
\renewcommand{\arraystretch}{1.4}
\begin{tabular}{|c|c|c|c|c|c|}
\hline
$x$ & $0$ & $1$ & $2$ & $3$ & $4$ \\
\hline
$f(x)$ & $5$ & $8$ & $11$ & $14$ & $17$ \\
\hline
\end{tabular}
\end{center}

\vspace{4pt}
(a)\quad Evaluate $f(3)$. \ans{$f(3) = 14$}

\vspace{4pt}
(b)\quad What is $f(2) - f(0)$? \ans{$11 - 5 = 6$}

\vspace{0.14in}
\textbf{2.}\quad The table below shows values of a function $g$.

\vspace{4pt}
\begin{center}
\renewcommand{\arraystretch}{1.4}
\begin{tabular}{|c|c|c|c|c|c|c|c|c|}
\hline
$x$ & $0$ & $1$ & $2$ & $3$ & $4$ & $5$ & $6$ & $7$ \\
\hline
$g(x)$ & $3$ & $7$ & $10$ & $7$ & $3$ & $7$ & $10$ & $7$ \\
\hline
\end{tabular}
\end{center}

\vspace{4pt}
(a)\quad Describe any pattern you notice in the output values of $g$.
\ans{The outputs repeat the sequence 3, 7, 10, 7 over and over, cycling every 4 units of input.}

\vspace{4pt}
(b)\quad Based on the pattern, what would you predict for $g(8)$?
\ans{$g(8) = 3$ (the pattern restarts at $x = 4$, so $x = 8$ is the same position in the next cycle)}

\vspace{0.14in}
\textbf{3.}\quad Using the table for $g$ above, compute the average rate of
change of $g$ from $x = 0$ to $x = 2$.

\vspace{4pt}
Average rate of change $= \dfrac{g(2) - g(0)}{2 - 0} = $ \ans{$\dfrac{10 - 3}{2} = 3.5$ per unit}

\vspace{0.14in}
\textbf{4.}\quad In your own words, what does the \emph{average rate of change}
of a function over an interval tell you?
\ans{It tells you, on average, how much the output changes for each one-unit increase in input over that interval --- the ``slope'' of the function between two points.}

\end{notesbox}

\end{document}
